\section{Пакетное вычисление статистик}

Рассмотрим наблюдения
\(
\Xv_i\in\R^{\dimd}
\), отклики
\(
Y_i\in\R
\), центры
\(
\xv_j
\), локализующие веса
\(
\weight_{ij}
\)
и направления
\(
\sph_{jr}\in\Sph_j
\), где
\(
i=1,\ldots,n
\),
\(
j=1,\ldots,\nJJJ
\),
\(
r=1,\ldots,\nSph
\).
В реализации направления хранятся в тензоре
\(
\Phi\in\R^{\nJJJ\times\nSph\times\dimd}
\), причём
\(
\Phi_{jr:}=\sph_{jr}^{\T}
\).

Для каждого центра определим массу окрестности, нормированные веса,
локальные средние наблюдений и откликов:
\begin{EQA}[c]
    \mu_j
    =
    \sum_{i=1}^{n}\weight_{ij},
    \qquad
    a_{ji}
    =
    \frac{\weight_{ij}}{\mu_j},
    \qquad
    \sum_{i=1}^{n}a_{ji}=1,
\end{EQA}
\begin{EQA}[c]
    \Mv_j
    =
    \sum_{i=1}^{n}a_{ji}\Xv_i,
    \qquad
    \bar Y_j
    =
    \sum_{i=1}^{n}a_{ji}Y_i.
\end{EQA}
Здесь предполагается
\(
\mu_j>0
\).
Вычисляемые статистики имеют вид
\begin{EQA}[c]
    \Ima_{j,\sph_{jr}}
    =
    \sum_{i=1}^{n}
    \weight_{ij}
    (Y_i-\bar Y_j)
    \left\langle
        \Xv_i-\Mv_j,
        \sph_{jr}
    \right\rangle,
\end{EQA}
\begin{EQA}[c]
    \Uv_{j,\sph_{jr}}
    =
    \sum_{i=1}^{n}
    \weight_{ij}
    (\Xv_i-\Mv_j)
    \left\langle
        \Xv_i-\Mv_j,
        \sph_{jr}
    \right\rangle.
\end{EQA}
Вместо отдельного прохода по каждому центру и каждому направлению код
одновременно обрабатывает пакет центров и все соответствующие им
направления.

\subsection{Разбиение по центрам}

Пусть
\(
\mathcal B=\{j_0,\ldots,j_0+B-1\}
\)
--- текущий пакет из
\(
B
\)
центров. Матрица его весов и матрица нормированных весов равны
\begin{EQA}[c]
    W_{\mathcal B}
    =
    (\weight_{ij})_{j\in\mathcal B,\,i=1,\ldots,n}
    \in\R^{B\times n},
    \qquad
    A_{\mathcal B}
    =
    \operatorname{diag}(\mu_{\mathcal B})^{-1}
    W_{\mathcal B}.
\end{EQA}
Если передана полная матрица весов
\(
W\in\R^{\nJJJ\times n}
\), она разбивается по строкам на пакеты размера
\texttt{batch\_size}. Можно также передать последовательность уже
сформированных блоков. Такие блоки должны без пропусков и перестановок
покрывать все центры; повторно они не разбиваются.

Сначала данные центрируются относительно общего среднего:
\begin{EQA}[c]
    \bar\Xv
    =
    \frac{1}{n}\sum_{i=1}^{n}\Xv_i,
    \qquad
    \Xv_i^c
    =
    \Xv_i-\bar\Xv.
\end{EQA}
Локальные средние центрированных данных вычисляются одним матричным
произведением:
\begin{EQA}[c]
    M_{\mathcal B}^c
    =
    A_{\mathcal B}X^c,
    \qquad
    \Mv_j^c
    =
    \sum_{i=1}^{n}a_{ji}\Xv_i^c.
\end{EQA}
После вычисления статистик восстанавливается исходное положение
локального среднего:
\begin{EQA}[c]
    \Mv_j
    =
    \Mv_j^c+\bar\Xv.
\end{EQA}
Глобальное центрирование уменьшает потерю точности, когда все
наблюдения имеют большой общий сдвиг.

Для всех центров и направлений текущего пакета одновременно строится
тензор проекций
\begin{EQA}[c]
    Q_{jri}
    =
    \left\langle
        \Xv_i^c-\Mv_j^c,
        \sph_{jr}
    \right\rangle,
    \qquad
    Q\in\R^{B\times\nSph\times n}.
\end{EQA}
В матричной форме это соответствует выражению
\begin{EQA}[c]
    Q
    =
    \Phi_{\mathcal B}(X^c)^{\T}
    -
    \left(
        \left\langle
            \Mv_j^c,
            \sph_{jr}
        \right\rangle
    \right)_{j,r,i}.
\end{EQA}
Последний член не зависит от
\(
i
\)
и распространяется вдоль третьей оси тензора.

В точной арифметике нормированное взвешенное среднее проекций равно
нулю. Из-за ошибок округления код вычисляет остаток
\begin{EQA}[c]
    e_{jr}
    =
    \sum_{i=1}^{n}a_{ji}Q_{jri}
\end{EQA}
и явно вычитает его:
\begin{EQA}[c]
    H_{jri}
    =
    a_{ji}(Q_{jri}-e_{jr}),
    \qquad
    s_{jr}
    =
    \sum_{i=1}^{n}H_{jri}.
\end{EQA}
После этого весь пакет статистик получается двумя матричными
произведениями:
\begin{EQA}[c]
    \Ima_{j,\sph_{jr}}
    =
    \mu_j
    \left(
        \sum_{i=1}^{n}H_{jri}Y_i
        -
        s_{jr}\bar Y_j
    \right),
\end{EQA}
\begin{EQA}[c]
    \Uv_{j,\sph_{jr}}
    =
    \mu_j
    \left(
        \sum_{i=1}^{n}H_{jri}\Xv_i^c
        -
        s_{jr}\Mv_j^c
    \right).
\end{EQA}
В точной арифметике
\(
e_{jr}=s_{jr}=0
\), поэтому эти выражения совпадают с исходными определениями.
Явное вычитание
\(
e_{jr}
\),
\(
s_{jr}\bar Y_j
\)
и
\(
s_{jr}\Mv_j^c
\)
подавляет накопившуюся ошибку нулевого момента.

\subsection{Плотное вычисление}

При плотном представлении весов в памяти одновременно находятся
основные массивы
\begin{EQA}[c]
    A_{\mathcal B}:B\times n,
    \qquad
    Q,H:B\times\nSph\times n,
    \qquad
    \Uv_{\mathcal B}:B\times\nSph\times\dimd.
\end{EQA}
Все направления одного центра вычисляются вместе, а независимые
матричные произведения выполняются сразу для всех центров пакета.
Основная вычислительная сложность равна
\begin{EQA}[c]
    O(B\nSph n\dimd),
\end{EQA}
а объём временной памяти ---
\begin{EQA}[c]
    O\left(
        B\nSph n
        +B\nSph\dimd
        +Bn
        +B\dimd
    \right).
\end{EQA}
Таким образом, параметр
\texttt{batch\_size}
ограничивает память, необходимую для тензоров
\(
Q
\)
и
\(
H
\), не меняя математическое определение статистик.

\subsection{Диагностические величины}

Помимо
\(
\Ima
\)
и
\(
\Uv
\), для каждого центра возвращаются масса
\(
\mu_j
\), локальное среднее
\(
\Mv_j
\)
и эффективное число наблюдений
\begin{EQA}[c]
    n_{\mathrm{eff},j}
    =
    \frac{1}{\sum_{i=1}^{n}a_{ji}^2}
    =
    \frac{
        \left(\sum_i\weight_{ij}\right)^2
    }{
        \sum_i\weight_{ij}^2
    }.
\end{EQA}
Чем меньше
\(
n_{\mathrm{eff},j}
\), тем сильнее масса окрестности сосредоточена на небольшом числе
наблюдений.

До вычитания остатка
\(
e_{jr}
\)
вычисляется относительная ошибка нулевого момента
\begin{EQA}[c]
    \eta_{jr}
    =
    \left\{
    \begin{array}{ll}
        \displaystyle
        \frac{
            |e_{jr}|
        }{
            \sum_i a_{ji}|Q_{jri}|
        },
        &
        \displaystyle
        \sum_i a_{ji}|Q_{jri}|\neq0,
        \\[4mm]
        0,
        &
        \text{иначе}.
    \end{array}
    \right.
\end{EQA}
Большое значение
\(
\eta_{jr}
\)
указывает на заметную потерю точности при центрировании проекций.

Итоговые массивы имеют формы
\begin{EQA}[c]
    \Ima:\nJJJ\times\nSph,
    \qquad
    \Uv:\nJJJ\times\nSph\times\dimd,
\end{EQA}
\begin{EQA}[c]
    \mu,n_{\mathrm{eff}}:\nJJJ,
    \qquad
    \Mv:\nJJJ\times\dimd,
    \qquad
    \eta:\nJJJ\times\nSph.
\end{EQA}
Полные массивы результатов сохраняются, а временные массивы текущего
пакета освобождаются перед обработкой следующего.
